\textbf{Protocol sensitivity is the finding.} The Phase-I collapse did not
reproduce under the clean equal-budget protocol: free ungated generation
already produced populations with 8.6\,pp oracle headroom and 31\% union
repair. Collapse is a failure mode of \emph{specific} generation pipelines ---
and it was invisible to every code-level check those pipelines had passed. The
practical corollary cuts against instinct: adding mechanism-verification did not
add routing value here, and its apparent cost was largely a population-size
artifact. Outcome-level auditing, not stronger code-level gating, is what
detects collapse. For sub-pp generation-side effects, the repeat-run noise
floor (14.5\% bare any-flip rate) should be reported alongside every CI --- it
exceeds the effects we estimate. Separately, the yield and headroom
decompositions show the gate changed \emph{which} candidates survive, not
\emph{whether} the population can answer: reliability and complementarity
need separate reporting.

\textbf{Limitations.} One domain family (text-to-SQL); Phase-I from a
discovery design with disclosed task overlap, a proxy judge, and a
post-trace human control; factorial inference conditional on the six sampled
builders; the cache-off arm replication (R2) is in progress; early Phase-I
generations lack per-attempt logs.
